\subsection{Reinforcement Learning}

Reinforcement Learning (RL) is a type of machine learning that focusses on an agent's ability to interact with its environment and make decisions. Unlike other types of learning methods, RL emphasises decision-making in potentially complex and uncertain environments. The goal of RL is to maximise rewards by learning from the consequences of actions, through trial and error. This approach has found applications in a range of fields such as robotics and finance, enabling the development of systems that improve their performance over time and adapt to new situations and make a series of decisions to achieve desired outcomes.

\subsubsection{Fundamentals and History}

Reinforcement Learning is a field of study that emerged from the combination of many other fields like psychology, computer science, and control theory, and it was inspired by the way humans and animals learn from their environment. RL involves an agent that makes decisions, receiving rewards or penalties based on its actions, aiming to maximise positive outcomes over time. This process is deeply rooted in the concept of operant conditioning, a type of learning discovered by psychologists like B. Skinner \cite{skinner_operant_1963}, where animal behaviour is shaped by its consequences. 

In the 1960s, Markov Decision Processes (MDPs) were introduced, which are frameworks that model decision making where the outcomes are partly random and partly under the control of the decision maker. This new field was the mathematical backbone of RL, as it allowed the development of algorithms that could learn optimal strategies in complex environments. One of the major breakthroughs in RL was the introduction of Q-Learning in the late 1980s, an algorithm capable of learning the value of actions in different states, helping agents make informed decisions without needing a model of the environment. This, along with Temporal Difference Learning, set the stage for the application of RL to a wide range of problems, from playing video games at a superhuman level to optimising financial trading strategies \cite{sutton_reinforcement_2018}.

\subsubsection{Key Concepts and Agent-Environment Interaction}

Reinforcement Learning is about how an agent interacts with its environment. This interaction involves states, actions and rewards. The state \( S_t \) represents the situation of the agent at that time instant. When the agent observes \( S_t \) it chooses an action \( A_t \). This action causes the environment to transition to a state \( S_{t+1} \), providing feedback in the form of a reward \( R_{t+1} \). Refer to Figure \ref{Figure:AgentEnvironment} for a visual representation of this cyclical interaction \cite{sutton_reinforcement_2018}. 

\begin{figure}[htb!]
\centering
\includegraphics[width=0.6\textwidth]{Images/AgentEnvironment.png}
\caption{Interaction between the agent and the environment in a Reinforcement Learning setting \cite{sutton_reinforcement_2018}.}
\label{Figure:AgentEnvironment}
\end{figure}

The main objective of the agent is to maximise its sum of rewards, also known as the return. The return is basically the expected gain that can be achieved starting from the state \( S_t \) as expressed by this formula \cite{sutton_reinforcement_2018}:

\begin{equation}
R_{t+1} + \gamma R_{t+2} + \gamma^2 R_{t+3} + \ldots = \sum_{k=0}^{\infty} \gamma^k R_{t+k+1}
\end{equation}

Here, \( \gamma \) represents the discount factor, indicating the degree to which future rewards are considered important relative to immediate ones. A policy \( \pi \) directs the agent's actions, which can be deterministic, assigning a specific action to each state, or stochastic, assigning a probability distribution over actions. Balancing between exploration, the act of trying new actions to discover their effects, and exploitation, the act of choosing actions that are known to yield high reward, is a fundamental challenge in RL. An important approach to this problem is the \(\varepsilon\)-greedy policy, which predominantly exploits the best-known action but occasionally explores random actions, thus ensuring both robustness and continual learning \cite{sutton_reinforcement_2018}.

The fundamental concept behind Reinforcement Learning lies in the Markov Decision Process (MDP) framework, defined by the tuple \( (S, A, P, R, \gamma) \). Here, \( S \) represents all possible states in the environment, \( A \) denotes the actions the agent can take, and \( P \) is the transition probability, indicating the likelihood of moving from state \( s_t \) to state \( s_{t+1} \) after taking action \( a \). \( R \) is the return received for this transition, and \( \gamma \) is the discount factor. MDPs utilise the Markov property, ensuring that the next state depends only on the current state and action, without need to consider any past states or actions, as shown in this formula \cite{sutton_reinforcement_2018}:

\begin{equation}
\label{Equation:MarkovProperty}
P(s_{t+1} | s_t, a_t, s_{t-1}, a_{t-1}, \ldots, s_0, a_0) = P(s_{t+1} | s_t, a_t)
\end{equation}

The reward function is very important as it captures the agents objectives. It offers a signal that directs the agent towards achieving the goals of a task, but it does not prescribe how to accomplish them. For example, in the trading context, rewards can mirror gains and losses, guiding the agent to adopt strategies that maximise profits \cite{sutton_reinforcement_2018}.

\subsubsection{Value Functions and Q-Functions}

The value function for a policy \( \pi \), denoted as \( V^\pi(s) \), represents the expected return when starting in state \( s \) and following policy \( \pi \) afterwards. It is formulated as follows \cite{sutton_reinforcement_2018}:

\begin{equation}
V^\pi(s) = \mathbb{E} \left[ \sum_{k=0}^{\infty} \gamma^k R_{t+k+1} \mid S_t = s, \pi \right]
\end{equation}

Here, \( \gamma \) is the discount factor, which weighs the importance of future rewards compared to immediate rewards, and \( \mathbb{E} \) denotes the expected value, reflecting the average of all possible outcomes weighted by their probabilities. The Q-function for a policy \( \pi \), denoted as \( Q^\pi(s, a) \), represents the expected return for taking an action \( a \) in state \( s \) and following policy \( \pi \) afterwards. It is formulated as follows \cite{sutton_reinforcement_2018}:

\begin{equation}
Q^\pi(s, a) = \mathbb{E} \left[ \sum_{k=0}^{\infty} \gamma^k R_{t+k+1} \mid S_t = s, A_t = a, \pi \right]
\end{equation}

The Bellman equations, named after Richard Bellman, who first formulated them, provide a recursive decomposition for the total cumulative reward an agent can expect to receive starting from a particular state or state-action pair. These equations are fundamental to dynamic programming techniques in RL, as they express the relationship between the value of the current state or state-action pairs as the expected values of subsequent states or state-action pairs. The recursive nature of the Bellman equations allows for the iterative improvement of the estimated values, essential for the convergence of RL algorithms. This is expressed as follows \cite{sutton_reinforcement_2018}:

\begin{equation}
V^\pi(s) = \sum_{a \in A} \pi(a|s) \sum_{s' \in S} P(s'|s,a) \left[ R(s,a,s') + \gamma V^\pi(s') \right]
\end{equation}

\begin{equation}
Q^\pi(s, a) = \sum_{s' \in S} P(s'|s,a) \left[ R(s,a,s') + \gamma \sum_{a' \in A} \pi(a'|s') Q^\pi(s', a') \right]
\end{equation}

\begin{equation}
\label{Equation:RelationshipValueQFunction}
V^\pi(s) = \sum_{a \in A} \pi(a|s) Q^\pi(s, a)
\end{equation}

The relationship pointed in Equation~\ref{Equation:RelationshipValueQFunction} indicates that the value of a state under policy \(\pi\) is the expected value of the Q-function over all possible actions, weighted by the policy's action probabilities. For optimal policies, we define the optimal value function \( V^*(s) \) and the optimal Q-function \( Q^*(s, a) \), which represent the maximum possible returns achievable from any state or state-action pair, respectively. Following the Bellman equation, it is possible to obtain \cite{sutton_reinforcement_2018}:

\begin{equation}
V^*(s) = \max_{a \in A} \sum_{s' \in S} P(s'|s,a) \left[ R(s,a,s') + \gamma V^*(s') \right]
\end{equation}

\begin{equation}
Q^*(s, a) = \sum_{s' \in S} P(s'|s,a) \left[ R(s,a,s') + \gamma \max_{a' \in A} Q^*(s', a') \right]
\end{equation}

\begin{equation}
V^*(s) = \max_{a \in A} Q^*(s, a)
\end{equation}

Through the iterative application of these equations, RL algorithms gradually approximate the value and Q-functions, guiding the agent towards the optimal policy that maximises the expected cumulative reward \cite{sutton_reinforcement_2018}.

\subsubsection{Methodological Variants and Learning Methods}

The field of reinforcement learning can be divided into categories based on the properties of the learning method: model-based versus model-free, on-policy versus off-policy strategies and value-based versus policy-based techniques. 

In model-based RL, algorithms create a model of the environment to simulate and predict the outcomes of actions. This allows the agent to plan by forecasting the states. This is possible when the environment is simple enough to be modelled. On the other hand, in model-free RL, algorithms directly learn a policy or value function from their interactions with the environment without explicitly modelling it. Model-free approaches are often chosen for environments where modelling is challenging or when an accurate model is not available. Within the realm of model-free RL, there are on-policy strategies, which learn about the policy that is currently used to make decisions, and off-policy strategies, which explore a different policy while aiming to learn an optimal one \cite{sutton_reinforcement_2018}.

Another important distinction lies between value-based (also designated as critic-only) and policy-based (also designated as actor-only) learning approaches. On the one hand, value-based methods, like Temporal Difference Learning (TD), Q-Learning and SARSA Learning, focus on estimating the value of states or state-action pairs. These algorithms derive policies implicitly from these values. SARSA Learning uses an on-policy strategy, while Q-Learning and TD Learning use an off-policy strategy. On the other hand, policy-based methods, such as REINFORCE Learning, directly optimise the policy itself, which is particularly useful when dealing with continuous action spaces. In particular, Actor-Critic learning combines elements of both policy-based and value-based approaches. It has a policy function (the actor) and a value function (the critic), which provides the stability and efficiency of value-based methods along with direct policy improvement. This collaboration between the actor and the critic results in improved policy updates, especially beneficial for complex environments \cite{sutton_reinforcement_2018}. Refer to Algorithm~\ref{Codes:ActorCritic} for a detailed implementation.

\begin{algorithm}[htb!]
\caption{Actor-Critic \cite{sutton_reinforcement_2018}.}
\label{Codes:ActorCritic}
\begin{algorithmic}
\REQUIRE $\pi(a \mid s; \theta)$: a differentiable policy, $Q(s, a; w)$: a differentiable $Q$-function, $N$: number of iterations and learning rates $\beta^\theta > 0$ and $\beta^w > 0$
\STATE Initialise policy parameter $\theta^{(0)}$ and $Q$-function parameter $w^{(0)}$
\STATE Initialise state $s$
\FOR{$n = 0, 1, \ldots, N - 1$}
    \STATE Sample $a \sim \pi(\cdot \mid s; \theta^{(n)})$
    \STATE Take action $a$, observe state $s'$ and reward $r$
    \STATE Sample $a' \sim \pi(\cdot \mid s'; \theta^{(n)})$
    \STATE Compute the TD error: $\delta \gets r + \gamma Q(s', a'; w^{(n)}) - Q(s, a; w^{(n)})$
    \STATE Update $w$: $w^{(n+1)} \gets w^{(n)} + \beta^w \delta \nabla_w Q(s, a; w^{(n)})$
    \STATE Update $\theta$: $\theta^{(n+1)} \gets \theta^{(n)} + \beta^\theta \delta \nabla_\theta \ln \pi(a \mid s; \theta^{(n)})$
    \STATE $s \gets s'$
\ENDFOR
\end{algorithmic}
\end{algorithm}

Although the Puterman theorem guarantees a deterministic stationary optimal policy for Markov Decision Processes (MDPs), such assurances do not extend to complex, non-stationary environments like financial markets \cite{sutton_reinforcement_2018}. Forex markets, in particular, are influenced by historical data in ways that standard state representations in MDPs do not account for. Consequently, RL in these markets relies on approximated models, utilising technical indicators and feature engineering to integrate essential historical information into the state representation in order to try to approximate the Markov property. This is why algorithms like the Actor-Critic learning are of special importance, because they act as a connection between reinforcement learning and neural networks, as we will see in the next section.